% ═══════════════════════════════════════════════════════════════════════════
\chapter{Objetivos}
% ═══════════════════════════════════════════════════════════════════════════

\section{Objetivo General}

Desarrollar un modelo de market making que integre control estocástico, métodos
semiparamétricos y machine learning para la estimación dinámica de intensidades de
órdenes, evaluando su desempeño empírico sobre datos de alta frecuencia de mercados
de criptoactivos.

\section{Objetivos Específicos}

\begin{enumerate}[label=\textbf{OE\arabic*.}]
  \item Analizar el modelo de \citet{AvellanedaStoikov2008} como caso base,
        identificando la forma exponencial en la distancia, los parámetros fijos y
        sus implicaciones sobre la política óptima de cotización.
  \item Proponer un modelo semiparamétrico para las intensidades de órdenes de la
        forma $\lambda_0(t) = \exp\bigl(g(X_t) + h(t)\bigr)$, especificando el vector
        de covariables $X_t$ y el componente temporal $h(t)$ a partir de propiedades
        documentadas de los datos de alta frecuencia.
  \item Implementar y comparar métodos de machine learning como Gradient Boosting
        y Random Forest para la estimación del componente no
        paramétrico $g(X_t)$, usando datos L2 históricos de BTC/USDT en Binance Spot.
  \item Derivar y resolver numéricamente la ecuación HJB bajo las intensidades
        semiparamétricas estimadas, obteniendo la política óptima de cotización
        data-driven.
  \item Evaluar el impacto del modelo propuesto frente a Avellaneda--Stoikov y a una
        política basada en Hawkes mediante las métricas de PnL neto, riesgo de
        inventario y robustez ante cambios de régimen de mercado.
\end{enumerate}

